\subsection{Гармоническая восьмёрка: минимальная степень слайса равна одиннадцати}

Положим
\[
 f=C_3=u^3-3u,\quad g=C_5=u^5-5u^3+5u,\quad
 h=C_7=u^7-7u^5+14u^3-7u,
\]
и рассмотрим над полем $k$ характеристики нуль
\[
 D_H=t\partial_u+f'(u)\partial_x+g'(u)\partial_y+h'(u)\partial_z.
\]
Здесь полная степень снова считается по всем пяти переменным $t,u,x,y,z$.

\paragraph{Теорема.}
Минимальная полная степень полиномиального слайса для $D_H$ равна $11$.
Доказательство ниже элементарно и не использует выпрямитель из репозитория.
Литературная новизна этого конкретного результата не заявляется.

\paragraph{Кубическая ретракция параметра.}
Тождества Чебышёва
\[
 C_3C_7-C_5^2=C_4-2,\qquad C_3C_4=C_7+u
\]
дают
\[
 R_H(x,y,z)=x^2z-xy^2+2x-z,\qquad R_H(f(u),g(u),h(u))=u.
\]
Степень $3$ этой ретракции минимальна: поскольку $f,g,h$ нечётны,
нечётная часть ограничения любого полинома степени не выше $2$
лежит в линейной оболочке $f,g,h$. Их различные степени $3,5,7$
не позволяют получить $u$. При этом $R_H$ не является координатой
объемлющего пространства: его градиент обращается в нуль в точках
$(1,0,-1)$ и $(-1,0,1)$.

\paragraph{Явный слайс степени одиннадцать.}
Положим $V=(f,g,h)-t(x,y,z)$. Точное полиномиальное частное
$\sigma_H=(u-R_H(V))/t$ имеет компактную форму
\begin{align*}
 \sigma_H={}&(2fh-g^2+2)x-2fg\,y+(f^2-1)z\\
 &+t\left(fy^2-2fxz+2gxy-hx^2\right)
 +t^2\left(x^2z-xy^2\right).
\end{align*}
После раскрытия скобок это $27$ ненулевых мономов полной степени $11$;
коэффициент при $u^{10}x$ равен $1$.
Из $D_H(V)=0$ сразу следует $D_H(\sigma_H)=1$;
тождество также проверено прямым точным дифференцированием.

\paragraph{Доказательство нижней границы.}
Используем вес $\operatorname{wt}(t)=1$, $\operatorname{wt}(u)=0$,
$\operatorname{wt}(X)=-1$, где $X=(x,y,z)$.
Дифференцирование повышает вес на единицу, поэтому гипотетический
слайс полной степени не выше $10$ можно заменить его компонентой веса $-1$:
\[
 s=\sum_{i=0}^{4}t^i s_i(u,X),\qquad
 \deg_Xs_i=i+1,\quad \deg_us_i\leq9-2i,
\]
причём каждое $s_i$ однородно по $X$.
Обозначим $D_0=f'\partial_x+g'\partial_y+h'\partial_z$.
Сравнивая коэффициенты в $D_Hs=1$, получаем
\[
 D_0s_0=1,\qquad \partial_us_i+D_0s_{i+1}=0\ (0\leq i<4),
 \qquad \partial_us_4=0.
\]

Последнее равенство делает $s_4$ постоянной формой степени $5$ от $X$.
В уравнении $\partial_us_3+D_0s_4=0$ первая часть имеет степень по $u$
не выше $2$. Сравнение коэффициентов при $u^6$ и $u^4$ даёт
$\partial_zs_4=\partial_ys_4=0$, откуда $s_4=cx^5$ и
\[
 s_3=-5cf x^4+K_4(X).
\]
В следующем уравнении $\deg_u\partial_us_2\leq4$.
Коэффициент при $u^6$ даёт $\partial_zK_4=0$;
коэффициент при $u^5$ равен $-60cx^3$, поскольку $D_0K_4$
содержит только чётные степени $u$. Значит, $c=0$ и $K_4=K_4(x,y)$.
Интегрирование теперь даёт
\[
 s_2=-f\,\partial_xK_4-g\,\partial_yK_4+K_3(X).
\]
В уравнении $\partial_us_1+D_0s_2=0$ первая часть имеет степень по $u$
не выше $6$. Коэффициенты при $u^9$ и $u^7$ последовательно дают
$\partial_{yy}K_4=\partial_{xy}K_4=0$. Так как $K_4$ однороден степени $4$,
он имеет вид $K_4=a x^4$. Следовательно,
\[
 s_2=-4af x^3+K_3(X),\qquad
 s_1=6af^2x^2-\gamma(u)\cdot\nabla K_3(X)+K_2(X).
\]
В уравнении $\partial_us_0+D_0s_1=0$ первая часть имеет степень по $u$
не выше $8$. Члены при $u^{13}$, $u^{11}$ и $u^9$ дают соответственно
\[
 \partial_{zz}K_3=0,\quad \partial_{yz}K_3=0,\quad
 \partial_{yy}K_3=-2\partial_{xz}K_3.
\]
Таким образом,
\[
 K_3=d x^3+e x^2y+k(x^2z-2xy^2).
\]
Ещё одно интегрирование даёт
\[
 s_0=-4af^3x+\tfrac12\gamma(u)^{\mathsf T}
      \operatorname{Hess}K_3(X)\gamma(u)
      -\gamma(u)\cdot\nabla K_2(X)+K_1(X),
\]
где $K_j$ --- однородная форма степени $j$ с постоянными коэффициентами.
По симметрии коэффициентов кубической формы равенство $D_0s_0=1$
теперь эквивалентно
\[
 \frac{d}{du}\left(K_1(\gamma)-K_2(\gamma)+K_3(\gamma)-af^4\right)=1.
\]
Так как $\gamma(0)=0$, выражение в скобках должно равняться $u$.
Его нечётная часть лежит в линейной оболочке
\[
 f,\quad g,\quad h,\quad f^3,\quad f^2g,\quad f^2h-2fg^2.
\]
Степени этих шести полиномов равны соответственно $3,5,7,9,11,13$;
старший коэффициент последнего равен $-1$.
Последовательное сравнение старших коэффициентов исключает возможность
получить из их линейной комбинации полином $u$. Противоречие доказывает
нижнюю границу и теорему.

\paragraph{Независимый точный сертификат.}
Помимо аналитического доказательства, построен рациональный функционал
извлечения коэффициентов со $117$ ненулевыми весами, для которого
$L(1)=1$ и $L(D_Hm)=0$ на всех $\binom{15}{5}=3003$ мономах полной
степени не выше $10$. Все веса находятся в файле
\texttt{no\_harmonic\_slice\_degree10\_certificate.json}.
Автономный \texttt{verify\_harmonic\_minimal\_slice.py} проверяет
каждое из этих равенств с помощью стандартной библиотеки Python и
точных рациональных чисел; он также проверяет компактную формулу
$\sigma_H$ и все $27$ её коэффициентов.
Независимый второй проверяющий код дал тот же результат.
Объединение мономов всех образов $D_Hm$ имеет $6017$ элементов
и максимальную степень $15$.
